\chapter{Доказательство теоремы о сходимости}
\label{app:proof}

В приложении приведены вспомогательные леммы и полное доказательство
Теоремы~\ref{thm:sgem}. Используются обозначения
раздела~\ref{ch:theory}: $\theta = (\boldsymbol{\alpha}_{1:K}, \mathbf{R})
\in \Theta = \mathcal{A}^K \times (\Delta_K)^M$ при конечном
$\mathcal{A}$, $\varepsilon_t$ --- пошаговая ошибка суррогата (в шкале
потерь, см.~\eqref{eq:eps_t}), $N = \sum_m |\mathcal{C}_m|$, а
вариационная свободная энергия, параметризованная функцией потерь $u$,
есть
\[
\mathcal{F}_u(q, \theta) \;=\; \sum_{m=1}^{M} |\mathcal{C}_m|
\sum_{k=1}^{K} q_{mk}\, \log \frac{r_{mk}\, e^{-u(\boldsymbol{\alpha}_k, \mathbf{r}_k)}}{q_{mk}}.
\]
Обозначим $L_u(\theta) = \max_q \mathcal{F}_u(q, \theta)$ ---
взвешенное по размерам лог-правдоподобие, порождённое потерей $u$, и
сократим $L_{\mathrm{true}} = L_{u_{\mathrm{true}}}$,
$L^{(t)} = L_{u^{(t)}}$.

\section{Вспомогательные леммы}
\label{app:lemmas}

\begin{proposition}[Максимизация ELBO]
\label{app:prop-elbo}
Для всякой функции потерь
$u: \mathcal{A} \times [0, 1]^M \to \mathbb{R}_{\geq 0}$ и всякого
$\theta\in\Theta$ выполнено
$L_u(\theta) = \max_{q} \mathcal{F}_u(q, \theta)$, причём максимум
достигается в
\[
q^{\star}_{mk}(\theta) \;=\; \frac{r_{mk}\,e^{-u(\boldsymbol{\alpha}_k, \mathbf{r}_k)}}
{\sum_{j} r_{mj}\,e^{-u(\boldsymbol{\alpha}_j, \mathbf{r}_j)}}.
\]
\end{proposition}

\begin{proof}
Зафиксируем $\theta$ и $m$. Множитель $|\mathcal{C}_m|$ положителен и
выносится из внутренней максимизации по $q_{m,:}$, не влияя на оптимум.
Обозначив $a_{mk} = r_{mk}\, e^{-u(\boldsymbol{\alpha}_k, \mathbf{r}_k)}$
и используя $\sum_k q_{mk} = 1$, внутреннее выражение равно
$-\mathrm{KL}(q_{m,:} \,\|\, a_{m,:}/\!\sum_j a_{mj}) + \log\sum_j a_{mj}$.
KL-дивергенция неотрицательна и обращается в нуль в точности при
$q_{mk} = a_{mk}/\sum_j a_{mj}$, что даёт утверждение. Суммирование по
$m$ с весами $|\mathcal{C}_m|$ восстанавливает
$L_u(\theta) = \sum_m |\mathcal{C}_m|\,\log\sum_k a_{mk}$.
\end{proof}

Следствием является разложение Нила--Хинтона
$L_u(\theta) = \mathcal{F}_u(q, \theta) + \sum_m |\mathcal{C}_m|\,
\mathrm{KL}(q_{m,:}\|q^{\star}_{m,:}(\theta))$, откуда
$L_u(\theta) \geq \mathcal{F}_u(q,\theta)$ с равенством тогда и только
тогда, когда $q = q^{\star}(\theta)$.

\begin{lemma}[Равномерное отклонение ELBO]
\label{app:lem-elbo}
При Предположении~\ref{ass:bounded} для всех $q$, $\theta$, $t$
выполнено
$|\mathcal{F}_{u_{\mathrm{true}}}(q,\theta) - \mathcal{F}_{u^{(t)}}(q,\theta)|
\leq N\,\varepsilon_t$.
\end{lemma}

\begin{proof}
Слагаемые $\log r_{mk}$ и $-\log q_{mk}$ сокращаются между двумя
ELBO, оставляя
$\mathcal{F}_{u_{\mathrm{true}}} - \mathcal{F}_{u^{(t)}}
= \sum_m |\mathcal{C}_m| \sum_k q_{mk}
[u^{(t)}(\boldsymbol{\alpha}_k, \mathbf{r}_k) - u_{\mathrm{true}}(\boldsymbol{\alpha}_k, \mathbf{r}_k)]$.
Оценивая скобку через $\varepsilon_t$ и используя $\sum_k q_{mk} = 1$,
получаем
$|\mathcal{F}_{u_{\mathrm{true}}} - \mathcal{F}_{u^{(t)}}|
\leq \varepsilon_t \sum_m |\mathcal{C}_m| = N\,\varepsilon_t$.
\end{proof}

\begin{corollary}
\label{app:cor-lik}
Для всякого $\theta$ выполнено
$|L_{\mathrm{true}}(\theta) - L^{(t)}(\theta)| \leq N\,\varepsilon_t$.
\end{corollary}

\begin{proof}
Применяя Лемму~\ref{app:lem-elbo} к максимумам двух ELBO и используя
$|\max_q f - \max_q g| \leq \sup_q |f - g|$.
\end{proof}

Для анализа предельных точек потребуются ещё три оценки --- на
softmax-отображение, на ответственности и на $Q$-функцию.

\begin{lemma}[Устойчивость softmax]
\label{app:lem-softmax}
Пусть $a_k, \tilde a_k > 0$ удовлетворяют
$e^{-\varepsilon} a_k \leq \tilde a_k \leq e^{\varepsilon} a_k$. Тогда
нормированные векторы $p_k = a_k/\sum_j a_j$ и
$\tilde p_k = \tilde a_k/\sum_j \tilde a_j$ удовлетворяют
$|\tilde p_k - p_k| \leq (e^{2\varepsilon} - 1)\,p_k \leq 3\varepsilon\,p_k$
при $\varepsilon \leq 1$.
\end{lemma}

\begin{proof}
$\tilde p_k \leq a_k e^{\varepsilon}/\sum_j a_j e^{-\varepsilon}
= p_k\,e^{2\varepsilon}$ и аналогично
$\tilde p_k \geq p_k\,e^{-2\varepsilon}$.
\end{proof}

\begin{lemma}[Отклонение ответственностей]
\label{app:lem-resp}
Пусть $w_{mk}(\theta) = q^{\star}_{mk}(\theta)$ ---
ответственности при истинном суррогате
(Утверждение~\ref{app:prop-elbo} при $u = u_{\mathrm{true}}$), а
$q^{(t)}_{mk}$ --- ответственности E-шага на итерации $t$. При
$\varepsilon_t \leq 1$ выполнено
$|q^{(t)}_{mk} - w_{mk}(\theta^{(t)})| \leq 3\varepsilon_t\, w_{mk}(\theta^{(t)})$.
\end{lemma}

\begin{proof}
Применяем Лемму~\ref{app:lem-softmax} с
$a_{mk} = r^{(t)}_{mk}\,e^{-u_{\mathrm{true}}(\boldsymbol{\alpha}^{(t)}_k, \mathbf{r}^{(t)}_k)}$
и $\tilde a_{mk}$, полученным заменой $u_{\mathrm{true}}$ на $u^{(t)}$.
Отношение $\tilde a_{mk}/a_{mk} = e^{u_{\mathrm{true}} - u^{(t)}}$
удовлетворяет
$e^{-\varepsilon_t} \leq \tilde a_{mk}/a_{mk} \leq e^{\varepsilon_t}$
ввиду~\eqref{eq:eps_t}, так что лемма применима при
$\varepsilon = \varepsilon_t$.
\end{proof}

\begin{lemma}[Отклонение $Q$-функции]
\label{app:lem-qdev}
Определим суррогатную $Q$-функцию
$\tilde Q^{(t)}(\theta;\theta') = \sum_{m} |\mathcal{C}_m|
\sum_{k} q^{(t)}_{mk}(\theta')[\log r_{mk} - u^{(t)}(\boldsymbol{\alpha}_k, \mathbf{r}_k)]$.
При Предположении~\ref{ass:bounded} существует постоянная
$C_Q = C_Q(N, K, U_{\max})$ такая, что для всех $\theta, \theta'$ и
всех $t$ с $\varepsilon_t \leq 1$ выполнено
$|\tilde Q^{(t)}(\theta;\theta') - Q_{\mathrm{true}}(\theta;\theta')|
\leq C_Q\,\varepsilon_t$.
\end{lemma}

\begin{proof}
Прибавим и вычтем
$\sum_m |\mathcal{C}_m| \sum_k q^{(t)}_{mk}[\log r_{mk} - u_{\mathrm{true}}(\boldsymbol{\alpha}_k, \mathbf{r}_k)]$.
Первая часть (член ошибки суррогата) есть
$\sum_m |\mathcal{C}_m| \sum_k q^{(t)}_{mk}(u_{\mathrm{true}} - u^{(t)})$
и ограничена $N\varepsilon_t$. Остаток распадается на
$\sum_m |\mathcal{C}_m| \sum_k (q^{(t)}_{mk} - w_{mk})(-u_{\mathrm{true}})$
и
$\sum_m |\mathcal{C}_m| \sum_k (q^{(t)}_{mk} - w_{mk})\log r_{mk}$.
Первое ограничено $3NK\,U_{\max}\,\varepsilon_t$ по
Лемме~\ref{app:lem-resp} (так как $|u_{\mathrm{true}}| \leq U_{\max}$).
Для второго, используя $|x\log x|\leq 1/e$ на $[0,1]$ вместе с оценкой
устойчивости softmax (Лемма~\ref{app:lem-softmax}), получаем оценку
$3N\varepsilon_t/e$. Суммируя три вклада, приходим к утверждению с
$C_Q = N(1 + 3K\,U_{\max} + 3/e)$.
\end{proof}

\section{Основное доказательство}
\label{app:mainproof}

\begin{proposition}[Квазимонотонность $L_{\mathrm{true}}$]
\label{app:prop-qm}
При Предположениях~\ref{ass:bounded}--\ref{ass:gem} для всякого $t$
выполнено
\begin{equation}
\label{eq:qm-app}
L_{\mathrm{true}}(\theta^{(t+1)}) \;\geq\;
L_{\mathrm{true}}(\theta^{(t)}) \;-\; 2N\,\varepsilon_t.
\end{equation}
\end{proposition}

\begin{proof}
Свяжем пять неравенств. Во-первых,
$L_{\mathrm{true}}(\theta^{(t+1)}) \geq
\mathcal{F}_{u_{\mathrm{true}}}(q^{(t)}, \theta^{(t+1)})$, так как
$L_u(\theta) = \max_q \mathcal{F}_u(q,\theta)$. Во-вторых, по
Лемме~\ref{app:lem-elbo},
$\mathcal{F}_{u_{\mathrm{true}}}(q^{(t)}, \theta^{(t+1)})
\geq \mathcal{F}_{u^{(t)}}(q^{(t)}, \theta^{(t+1)}) - N\varepsilon_t$.
В-третьих, обобщённый M-шаг (Предположение~\ref{ass:gem}) даёт
$\mathcal{F}_{u^{(t)}}(q^{(t)}, \theta^{(t+1)})
\geq \mathcal{F}_{u^{(t)}}(q^{(t)}, \theta^{(t)})$. В-четвёртых, E-шаг
делает последнее равным $L^{(t)}(\theta^{(t)})$. В-пятых,
Следствие~\ref{app:cor-lik} даёт
$L^{(t)}(\theta^{(t)}) \geq L_{\mathrm{true}}(\theta^{(t)}) - N\varepsilon_t$.
Объединяя пять шагов, получаем~\eqref{eq:qm-app}.
\end{proof}

\textbf{Часть (a): сходимость правдоподобия.} Положим
$L_{\max} = \sup_{\theta\in\Theta} L_{\mathrm{true}}(\theta)$ (величина
конечна ввиду компактности $\Theta$ --- Предположение~\ref{ass:cont} и
конечность $\mathcal{A}$). Определим
$V_t = L_{\max} - L_{\mathrm{true}}(\theta^{(t)}) \geq 0$. Из
Утверждения~\ref{app:prop-qm} следует
$V_{t+1} \leq V_t + 2N\varepsilon_t$, причём
$\sum_{t} 2N\varepsilon_t < \infty$ по Предположению~\ref{ass:summable}.
По теореме Роббинса--Зигмунда~\cite{robbins1971convergence} (с $a_t = 0$,
$b_t = 2N\varepsilon_t$) последовательность $\{V_t\}$ сходится, а значит,
сходится и $\{L_{\mathrm{true}}(\theta^{(t)})\}$. \qed

\textbf{Часть (b): стационарность предельных точек.}

\emph{Шаг 1: существование.} $\Theta$ компактно, поэтому у
$\{\theta^{(t)}\}$ есть хотя бы одна предельная точка
$\theta^{\star} = \lim_{j\to\infty} \theta^{(t_j)}$.

\emph{Шаг 2: предположение от противного.} Пусть $\theta^{\star}$ не
стационарна: найдутся $\theta'\in\Theta$ и $\delta > 0$ с
$Q_{\mathrm{true}}(\theta';\theta^{\star}) -
Q_{\mathrm{true}}(\theta^{\star};\theta^{\star}) = \delta > 0$.

\emph{Шаг 3: непрерывность.} Отображение
$\theta'\mapsto Q_{\mathrm{true}}(\theta;\theta')$ непрерывно на
$\Theta$: $w_{mk}(\theta')$ непрерывно, так как $u_{\mathrm{true}}$
непрерывно по $\mathbf{r}$ (Предположение~\ref{ass:cont}), а softmax
непрерывен; произведение $w_{mk}\log r_{mk}$ непрерывно продолжается в
$r_{mk} = 0$, так как $w_{mk} = O(r_{mk})$, а $r\mapsto r\log r$
непрерывно в нуле. Поэтому при достаточно больших $j$
$Q_{\mathrm{true}}(\theta';\theta^{(t_j)}) -
Q_{\mathrm{true}}(\theta^{(t_j)};\theta^{(t_j)}) \geq \delta/2$.

\emph{Шаг 4: переход к $\tilde Q$.} По Лемме~\ref{app:lem-qdev},
$\tilde Q^{(t_j)}(\theta';\theta^{(t_j)}) -
\tilde Q^{(t_j)}(\theta^{(t_j)};\theta^{(t_j)}) \geq
\delta/2 - 2C_Q\,\varepsilon_{t_j} \geq \delta/4$
при достаточно больших $j$, так как $\varepsilon_{t_j}\to 0$.

\emph{Шаг 5: подъём на M-шаге.} Поскольку M-шаг перебирает множество
кандидатов, содержащее $\theta'$ (ввиду конечности $\mathcal{A}$ и
свойств замыкания обновления по $\mathbf{R}$),
$\tilde Q^{(t_j)}(\theta^{(t_j+1)};\theta^{(t_j)}) \geq
\tilde Q^{(t_j)}(\theta';\theta^{(t_j)})$. Переходя обратно через
Лемму~\ref{app:lem-qdev}, получаем
$\Delta_{t_j} := Q_{\mathrm{true}}(\theta^{(t_j+1)};\theta^{(t_j)}) -
Q_{\mathrm{true}}(\theta^{(t_j)};\theta^{(t_j)}) \geq
\delta/4 - 2C_Q\,\varepsilon_{t_j} \geq \delta/8$
при достаточно больших $j$.

\emph{Шаг 6: противоречие через суммируемость.} Неравенство Гиббса
$L_{\mathrm{true}}(\theta^{(t+1)}) - L_{\mathrm{true}}(\theta^{(t)})
\geq \Delta_t$ вместе с Утверждением~\ref{app:prop-qm} даёт
$\Delta_t \geq -2N\varepsilon_t$. Тогда неотрицательная
последовательность $\Delta_t + 2N\varepsilon_t$ удовлетворяет
$\sum_t (\Delta_t + 2N\varepsilon_t) < \infty$ (Роббинс--Зигмунд для
$V_t$), и потому $\sum_t \Delta_t < \infty$. Но Шаг 5 даёт
$\Delta_{t_j} \geq \delta/8$ для бесконечного числа $j$, откуда
$\sum_t \Delta_t = +\infty$ --- противоречие. Следовательно,
$\theta^{\star}$ стационарна. \qed

\section{Скорость сходимости}
\label{app:rate}

Побочным следствием доказательства части~(a) является явная оценка
скорости
\[
L_{\mathrm{true}}(\theta^{\star}) - L_{\mathrm{true}}(\theta^{(T)})
\;\leq\; 2N \sum_{t = T}^{\infty} \varepsilon_t.
\]
Если $\varepsilon_t = O(t^{-p})$ при $p > 1$, то
$L_{\mathrm{true}}(\theta^{(T)}) \to L_{\mathrm{true}}(\theta^{\star})$
со скоростью $O(T^{1 - p})$, что согласуется со скоростями анализа
стохастического EM~\cite{cappe2009online}.
